\documentclass[11pt]{article}
\usepackage[margin=1in]{geometry}
\usepackage[T1]{fontenc}
\usepackage[utf8]{inputenc}
\usepackage{lmodern}
\usepackage{graphicx}
\usepackage{booktabs}
\usepackage{tabularx}
\usepackage{array}
\usepackage{hyperref}
\usepackage{ragged2e}

\newcolumntype{Y}{>{\RaggedRight\arraybackslash}X}
\renewcommand{\arraystretch}{1.15}

\title{Supplementary Material for\\Optimization Against AI Slop}
\author{PromptLab / STAT 4830 Repository Synthesis}
\date{\today}

\begin{document}
\maketitle

\section*{S1. Figure Provenance and Asset Inventory}

All figures used in the revised report are either derived directly from repository output files or reconstructed as clean academic diagrams from repository code paths and presentation notes. Slide screenshots are intentionally excluded from the report and from the recommended figure set below.

\begin{table}[h]
\centering
\small
\begin{tabularx}{\linewidth}{l Y Y}
\toprule
Asset & Source basis & Recommended use \\
\midrule
\texttt{tournament\_phase2\_tradeoff.png} & Replotted from \texttt{tournament/results/phase2/leaderboard.csv}. & Main paper figure for the tournament section. Shows the tradeoff between score improvement and fidelity. \\
\texttt{editlens\_reference\_bands.png} & Reconstructed from the project's EditLens calibration notes. & Main paper figure for the RL setup section. Explains the score geometry. \\
\texttt{rl\_progression\_clean.png} & Replotted from \texttt{outputs/c2\_eval.jsonl}. & Main paper figure for the RL results section. Summarizes both score trajectory and label-bucket transfer. \\
\texttt{kl\_ablation\_clean.png} & Reconstructed from \texttt{outputs/c2\_eval.jsonl}, \texttt{outputs/c4a\_eval.jsonl}, and \texttt{outputs/c4c\_eval.jsonl}. & Main paper figure for reward-hacking analysis. \\
\texttt{promptlab\_v2\_pipeline.png} & Reconstructed from code paths in \texttt{src/hill\_climb/prompt\_opt/} and \texttt{src/hill\_climb/slop\_gen/}, plus presentation notes. & Main paper figure for the Promptlab-v2 method section. \\
\texttt{promptlab\_v2\_results.png} & Reconstructed from the Promptlab-v2 presentation notes. & Main paper figure for the Promptlab-v2 validation result. \\
\bottomrule
\end{tabularx}
\end{table}

Recommended compiled-paper subset:

\begin{itemize}
\item Keep all six figures above in the main paper if space allows.
\item If page pressure becomes severe, move \texttt{promptlab\_v2\_pipeline.png} to the supplement before cutting any RL figure.
\end{itemize}

\section*{S2. Extended Quantitative Notes}

\subsection*{S2.1 RL checkpoint trajectory (\texttt{outputs/c2\_eval.jsonl})}

\begin{table}[h]
\centering
\small
\resizebox{\linewidth}{!}{
\begin{tabular}{rrrrrr}
\toprule
Step & Mean EditLens & LABEL\_0 (\%) & LABEL\_1 (\%) & LABEL\_3 (\%) & KL \\
\midrule
0 & 0.9990 & 0 & 0 & 100 & 0.0000 \\
50 & 0.7677 & 1 & 29 & 70 & 0.3972 \\
100 & 0.7664 & 0 & 32 & 68 & 0.4215 \\
150 & 0.7535 & 0 & 36 & 64 & 0.4567 \\
200 & 0.7246 & 3 & 39 & 58 & 0.4814 \\
250 & 0.6819 & 4 & 44 & 52 & 0.5133 \\
300 & 0.7026 & 5 & 35 & 60 & 0.5495 \\
350 & 0.7178 & 5 & 33 & 62 & 0.5513 \\
400 & 0.6749 & 11 & 31 & 58 & 0.5804 \\
450 & 0.6850 & 9 & 34 & 56 & 0.5447 \\
500 & 0.6700 & 11 & 34 & 55 & 0.5183 \\
\bottomrule
\end{tabular}
}
\end{table}

Interpretation:

\begin{itemize}
\item The main gain happens early: by step 50 the model has already moved roughly 23 percentage points in mean EditLens.
\item The training curve is not monotone after step 250, which is consistent with a high-variance RL regime.
\item Later gains are incremental and partly offset by instability.
\end{itemize}

\subsection*{S2.2 KL-penalty ablation}

\begin{table}[h]
\centering
\small
\resizebox{\linewidth}{!}{
\begin{tabular}{lrrl}
\toprule
Run & Mean EditLens at step 500 & KL divergence at step 500 & Interpretation \\
\midrule
C2 (\texttt{KL=0.1}) & 0.6700 & 0.5183 & best trustworthy operating point \\
C4-a (\texttt{KL=0.0}) & 0.6325 & 0.6371 & slightly better detector score, somewhat less stable \\
C4-c (\texttt{KL=0.5}) & 0.4307 & 3.7347 & apparent improvement is not trustworthy; strong reward-hacking signature \\
\bottomrule
\end{tabular}
}
\end{table}

\subsection*{S2.3 Tournament phase-1 summary}

\begin{table}[h]
\centering
\small
\begin{tabular}{lrrr}
\toprule
Run & Rows & Mean delta EditLens & Median length ratio \\
\midrule
M1 & 500 & -0.00639 & 0.7471 \\
M2 & 500 & -0.00674 & 0.7471 \\
M3 & 500 & -0.00644 & 0.7471 \\
M4 & 500 & -0.01251 & 0.7519 \\
\bottomrule
\end{tabular}
\end{table}

\subsection*{S2.4 Tournament phase-2 summary}

\begin{table}[h]
\centering
\small
\resizebox{\linewidth}{!}{
\begin{tabular}{lrrrr}
\toprule
Run & Pass rate & Mean delta EditLens & Mean similarity & Median length ratio \\
\midrule
M4 & 0.160 & +0.00035 & 0.8486 & 0.5981 \\
M7 & 0.128 & -0.02908 & 0.8981 & 0.5851 \\
M3 & 0.104 & -0.06519 & 0.9628 & 0.5776 \\
M5 & 0.096 & -0.06308 & 0.9663 & 0.5776 \\
M2 & 0.092 & -0.06510 & 0.9610 & 0.5795 \\
M1 & 0.080 & -0.06515 & 0.9640 & 0.5783 \\
M6 & 0.076 & -0.06465 & 0.9619 & 0.5784 \\
\bottomrule
\end{tabular}
}
\end{table}

Interpretation:

\begin{itemize}
\item \texttt{M4} is nominally the leader, but the effect size is tiny and accompanied by poor length preservation.
\item The shared harness makes it difficult to claim a meaningful winner.
\end{itemize}

\subsection*{S2.5 Multi-detector summary}

The 20-essay multi-detector check reports mean EditLens \texttt{0.9900 -> 0.8171} and a perplexity-based statistic \texttt{10.94 -> 57.10}. This should not be read as straightforward generalization. The rise in the second statistic is dominated by a few degenerate outputs with extreme values, including one case above 800.

\section*{S3. Qualitative Failure Modes}

\subsection*{S3.1 RL branch}

\begin{itemize}
\item Essay 5 (\texttt{outputs/e3\_qualitative.txt}): the model refuses to rewrite violent content and instead outputs a safety refusal. Detector score improves, but the rewrite task is not completed.
\item Essay 6 (\texttt{outputs/e3\_qualitative\_v2.txt}): strongest clean success case. The model preserves the detective-story setup while making denser stylistic edits; score improves from \texttt{0.9995} to \texttt{0.4764}.
\item Essay 7 (\texttt{outputs/e3\_qualitative\_v2.txt}): the model falls into broken token sequences and pseudo-words while still driving the detector score down sharply.
\item Essay 9 (\texttt{outputs/e3\_qualitative.txt} and \texttt{outputs/e3\_qualitative\_v2.txt}): the model begins a plausible rewrite and then spills into explanation or malformed continuation, illustrating partial reward hacking rather than stable editing.
\end{itemize}

\subsection*{S3.2 Promptlab-v2 branch}

The Promptlab-v2 notes describe an example-essay comparison in which the strongest model (\texttt{v3}) outperforms both the baseline and the prompt-to-prompt alternatives, but does so by producing visibly shorter prose. The repository's own interpretation is the right one: this is a T5-base artifact caused by summarization pretraining, not evidence that the overall search-and-distill idea is invalid.

\section*{S4. Suggested Supplementary Tables for Future Expansion}

If this report is later expanded into a submission-length draft, the following tables would strengthen the supplement:

\begin{enumerate}
\item a run-by-run transcription of the Promptlab-v2 ablation table described in the presentation notes;
\item a small table of the exact LoRA targets, learning rates, and stopping criteria across \texttt{M1}--\texttt{M7};
\item a detector-regime table that explicitly separates local slop metrics, Pangram RoBERTa, and Pangram 3B;
\item a human-evaluation placeholder table, since the strongest remaining limitation across all branches is the absence of human judgment data.
\end{enumerate}

\end{document}
